\documentclass[10pt,twocolumn,letterpaper]{article}
\usepackage[rebuttal]{cvpr}

% Include other packages here, before hyperref.
\usepackage{graphicx}
\usepackage{amsmath}
\usepackage{amssymb}
\usepackage{booktabs}

% Import additional packages in the preamble file, before hyperref
\input{preamble}

% If you comment hyperref and then uncomment it, you should delete
% egpaper.aux before re-running latex.  (Or just hit 'q' on the first latex
% run, let it finish, and you should be clear).
\definecolor{cvprblue}{rgb}{0.21,0.49,0.74}
\usepackage[pagebackref,breaklinks,colorlinks,allcolors=cvprblue]{hyperref}

% If you wish to avoid re-using figure, table, and equation numbers from
% the main paper, please uncomment the following and change the numbers
% appropriately.
%\setcounter{figure}{2}
%\setcounter{table}{1}
%\setcounter{equation}{2}

% If you wish to avoid re-using reference numbers from the main paper,
% please uncomment the following and change the counter value to the
% number of references you have in the main paper (here, 100).
%\makeatletter
%\apptocmd{\thebibliography}{\global\c@NAT@ctr 100\relax}{}{}
%\makeatother

%%%%%%%%% PAPER ID  - PLEASE UPDATE
\def\paperID{*****} % *** Enter the Paper ID here
\def\confName{CVPR}
\def\confYear{2026}

\begin{document}

%%%%%%%%% TITLE - PLEASE UPDATE
\title{\LaTeX\ CoMa: Contextual Massing Generation with Vision-Language Models}  % **** Enter the paper title here

\maketitle
\thispagestyle{empty}
\appendix

%%%%%%%%% BODY TEXT - ENTER YOUR RESPONSE BELOW
\section{Introduction}

Thank you for the constructive reviews and for acknowledging the value of the CoMa-20K dataset. We appreciate the opportunity to clarify the scope and contributions of our work in response to your concerns.

We fully acknowledge the limitation that CoMa-20K is derived from a single city, Melbourne. Creating a large-scale, multimodal dataset of this kind, which aligns geometric, textual, and visual urban data from real-world sources, is a significant undertaking. Our primary goal with this initial release is to establish a foundational benchmark to catalyze research in data-driven massing generation, a field currently held back by the scarcity of such resources. We present this dataset as a necessary and open starting point, with the explicit understanding that its expansion to cover diverse global urban contexts is a vital future direction for the community.

Regarding the request for non-VLM baselines, we agree that broader comparisons would strengthen the evaluation. Our focus was on establishing the dataset and providing initial benchmarks using a prevalent and accessible class of models. While other methods for geometric generation exist, they often address different task formulations, lack publicly available code, or require significant adaptation to fit our multimodal, context-conditioned framework. Benchmarking against these approaches is an important next step, and we will frame this clearly as a direction for future work enabled by the release of CoMa-20K.

We wish to reiterate that the core contribution of this paper is the dataset and benchmark itself, not a novel model architecture. A key motivation of experiments was to test whether fine-tuning could address the geometric reasoning issues known to plague zero-shot VLMs, which are often used in generative design due to a lack of training data. Our results show that even with dedicated training, standard VLM fine-tuning on tokenized geometry remains fundamentally problematic, leading to poor geometric adherence. We believe presenting this clear negative result is valuable for the community, as it definitively shows the limitations of a common approach and redirects focus toward the need for more robust geometric representations or hybrid methods.

We thank the reviewers for highlighting the apparent inconsistency between quantitative and qualitative results. This discrepancy is instructive. The zero-shot model achieves higher quantitative scores (e.g., Site IoU) by generating simplistic, often rectangular forms that roughly fit the site boundary and meet floor area targets. In contrast, the fine-tuned models attempt to produce more complex and architecturally interesting massings but suffer from geometric artifacts that severely penalize the same geometric metrics. The qualitative comparison thus reveals a trade-off: the fine-tuned models learn stylistic patterns but lack geometric robustness, while the zero-shot model defaults to safe, simplistic solutions.

In conclusion, we position this work squarely as a dataset and benchmark paper. Its significance lies in providing the first large-scale, multimodal resource to open a pathway for rigorous research in contextual massing generation. By sharing CoMa-20K with the community we aim to provide a solid foundation upon which future, more methodologically novel work can be built.

\end{document}
